% =====================================================================
% Apart Research hackathon submission template — LaTeX port
%
% Faithful reproduction of "Copia de Apart Research hackathon submission
% template" (.docx / .md / .pdf) in this same folder. Replace the
% [PLACEHOLDER] text and the italicized guidance paragraphs with your
% content, then delete the yellow guidance box before submitting, exactly
% as the original instructs.
% =====================================================================
\documentclass[11pt]{article}

% ---------- layout ----------
\usepackage[letterpaper,margin=1in]{geometry}
\usepackage[utf8]{inputenc}
\usepackage[T1]{fontenc}
\usepackage{parskip}          % blank line between paragraphs, no indent — matches the source doc
\usepackage{xcolor}
\usepackage{enumitem}
\usepackage{tcolorbox}
\usepackage{titlesec}
\usepackage[normalem]{ulem}   % \uline: underline that can break across lines (unlike \underline)
\usepackage{graphicx}
\usepackage{booktabs}
\usepackage{longtable}
\usepackage[colorlinks=true,linkcolor=blue,urlcolor=blue,citecolor=blue]{hyperref}

% ---------- section headings: bold black serif, matches "1. Introduction" ----------
\titleformat{\section}{\normalfont\Large\bfseries}{\thesection.}{0.5em}{}
\titlespacing*{\section}{0pt}{1.6em}{0.8em}

% unnumbered gray subsection headings: matches "Limitations" / "Future Work"
\titleformat{\subsection}{\normalfont\large\bfseries\color{gray!70!black}}{}{0em}{}
\titlespacing*{\subsection}{0pt}{1.2em}{0.5em}

\setlist[enumerate,1]{leftmargin=1.8em}
\setlist[itemize,1]{leftmargin=1.8em}

\pagestyle{plain}

\begin{document}

% ===================== TITLE BLOCK =====================
\noindent\rule{\linewidth}{1.1pt}
\begin{center}
  \vspace{0.6em}
  {\LARGE\bfseries PROJECT TITLE\footnote{Research conducted at the
  \href{https://apartresearch.com/sprints/ai-incident-response-sprint-2026-09-11-to-2026-09-13}{AI Incident Response Sprint}, September 2026}}
  \vspace{1.0em}
\end{center}
\noindent\rule{\linewidth}{1.1pt}
\vspace{1.2em}

\begin{center}
\begin{tabular}{ccc}
Author name 1 & Author name 2 & Author name 3 \\
Affiliation   & Affiliation   & Affiliation   \\[0.9em]
Author name 4 & Author name 5 & Author name 6 \\
Affiliation   & Affiliation   & Affiliation   \\[0.9em]
\end{tabular}

\vspace{0.3em}
\textbf{With}\\
Apart Research
\end{center}

\vspace{1.2em}
\begin{center}
{\large\bfseries Abstract}
\end{center}
\begin{quote}
\itshape
Summarize your project in 150--250 words. A strong abstract lets a reviewer understand what
you did and why it matters without reading anything else. \uline{Make sure to cover: the
problem, your approach, key results, and the main takeaway.} Polish it last: the abstract
should reflect your final results, not your initial plan.
\end{quote}

\vspace{1.0em}
\begin{tcolorbox}[colback=yellow!12,colframe=black!70,boxrule=0.6pt,arc=0pt,left=8pt,right=8pt,top=6pt,bottom=6pt]
\textbf{How to use this template}: Replace the \textit{italicized guidance text} under each
section with your content. The section structure is strong guidance but not rigid. If your
project requires a different organization, feel free to adapt. \uline{Delete all guidance
text including this info box before submitting.} Make sure the project title and author
information above are up to date.

\vspace{0.6em}
\textbf{How your submission will be evaluated}: Your project will be judged on the quality of
this written report. To understand what we look for, see our
\href{https://apartresearch.notion.site/sprint-evaluation-rubric}{Evaluation Rubric}.

\vspace{0.6em}
\textbf{Recommended length}: 4 pages excluding references and appendix.\\
Rough guide: Intro \& Related Work 1p, Methods and Results: 2.5p, Discussion 0.5p.
\end{tcolorbox}

\vspace{1.4em}

% ===================== 1. INTRODUCTION =====================
\section{Introduction}

\textit{What problem are you addressing and why does it matter?
Connect it to your work: we want to know why your work is practically valuable.}

\textit{Provide enough background for readers to understand your work.
If relevant, briefly describe the threat model or failure mode you're addressing; reference
prior work that motivates why it is worth addressing, or explain it yourself.}

\textit{Aspire to clearly list your most important contributions that go beyond what exists
today.}\\
\textit{\textbf{Our main contributions are:}}

\begin{enumerate}
  \item \textit{[First contribution --- what new thing did you create, discover, or demonstrate?]}
  \item \textit{[Second contribution]}
  \item \textit{[Third contribution, if applicable]}
\end{enumerate}

% ===================== 2. RELATED WORK =====================
\section{Related Work}

\textit{What prior work is most similar, and how does your work differ?
Cite the most relevant papers, tools, or projects. Explain what gap your work addresses.}

\textit{Some questions which may help:}

\begin{enumerate}[label=\arabic*)]
  \item \textit{When and why would someone use your method over the existing state-of-the-art?}
  \item \textit{What information/insight does your method provide which we did not have before?}
\end{enumerate}

% ===================== 3. METHODS =====================
\section{Methods}

\textit{Describe your approach clearly enough that someone could replicate it. Include key
design choices and justify them where relevant (hint: the more you can back up your design
choices by referencing prior work, the better).}

\textit{What models/datasets/tools did you use and why? What were key parameters or design
decisions? What did you try that didn't work? Could someone reproduce your work from this
description?}

% ===================== 4. RESULTS =====================
\section{Results}

\textit{Present your main findings with appropriate evidence. Use figures and tables where
appropriate (we strongly encourage at least one figure, see tips below). Distinguish between
observations and interpretations.}

\textit{Argue why your claims are robust. E.g., if your approach ``performs better'' than
alternatives:}

\begin{enumerate}[label=\arabic*)]
  \item \textit{Do you have enough data? Is the difference statistically significant?}
  \item \textit{Is it robust? Or do small changes to your setup cause substantial changes to
  your results?}
\end{enumerate}

\begin{quote}
\textit{\textbf{Tips for figures and tables:}}
\begin{itemize}[label=--]
  \item \textit{Number all figures (Figure 1, Figure 2\ldots) and tables (Table 1, Table
  2\ldots)}
  \item \textit{Include descriptive captions that can be understood without the main text}
  \item \textit{Place figures/tables near where they're first referenced}
  \item \textit{\textbf{Ensure text in figures is legible!}}
\end{itemize}
\end{quote}

% ===================== 5. DISCUSSION AND LIMITATIONS =====================
\section{Discussion and Limitations}

\textit{Discuss the broader implications for AI safety.
What do your results mean? What trends do you notice and what might they indicate?}

\subsection*{Limitations}

\textit{What are the limitations of your work? What threat models or failure modes did you not
address? Be honest about constraints --- methodological limitations, scope limitations, or
aspects you couldn't fully address in the hackathon timeframe. Explicitly note the assumptions
you made, whether implicitly or explicitly, and how the interpretation of your results would
change if a given assumption did not hold.}

\subsection*{Future Work}

\textit{What are the natural next steps? How could this work be extended?}

% ===================== 6. CONCLUSION =====================
\section{Conclusion}

\textit{Briefly summarize your main findings and their implications (1--2 paragraphs).}

% ===================== CODE AND DATA =====================
\section*{Code and Data}

\textit{Include links if applicable. If your project doesn't involve code (e.g., policy
analysis) or if there are info-hazard considerations, note that here.}

\begin{itemize}[label=--]
  \item \textit{\textbf{Code repository}: [Link to GitHub/GitLab if applicable]}
  \item \textit{\textbf{Data/Datasets}: [Link if applicable]}
  \item \textit{\textbf{Other artifacts} (optional): [Demo link, video walkthrough, Hugging Face
  Space, etc.]}
\end{itemize}

% ===================== AUTHOR CONTRIBUTIONS =====================
\section*{Author Contributions (optional)}

\textit{[e.g., ``A.B. led the project and designed experiments. C.D. implemented the code. All
authors contributed to writing and reviewed the final manuscript.'']}

% ===================== REFERENCES =====================
\section*{References}

\textit{Use a consistent citation format. Include: Author(s), Year, Title, Venue/Publisher, and
URL or DOI where available.}

\begin{enumerate}
  \item \textit{[Reference 1]}
  \item \textit{[Reference 2]}
  \item \textit{\ldots}
\end{enumerate}

% ===================== APPENDIX =====================
\section*{Appendix \textit{(optional)}}

\textit{Supplementary material such as additional figures, detailed methodology, prompts used,
extended results, etc.}

% ===================== LLM USAGE STATEMENT =====================
\section*{LLM Usage Statement}

\textit{If you used LLM assistance in developing your project or writing this report, briefly
note how. Ensure all claims and results have been verified.}

\textit{NOTE: We strongly encourage that the final version of the submission is primarily
written by your team.}

\textit{[e.g., ``We used Claude to brainstorm approaches and help draft sections. All results
and claims were independently verified.'']}

\end{document}
